\documentclass[UTF8,11pt]{ctexart}
\usepackage[a4paper,margin=24mm]{geometry}
\usepackage{amsmath,amssymb,amsthm,mathtools}
\usepackage[all]{xy}
\usepackage{xurl,hyperref,enumitem}
\hypersetup{hidelinks,breaklinks=true}
\setlength{\emergencystretch}{3em}
\newtheorem*{theorem}{Theorem}
\newtheorem*{lemma}{Lemma}
\newtheorem*{proposition}{Proposition}
\newtheorem*{corollary}{Corollary}
\theoremstyle{definition}
\newtheorem*{definition}{Definition}
\newtheorem*{remark}{Remark}
\DeclareMathOperator{\Hom}{Hom}
\DeclareMathOperator{\Ext}{Ext}
\DeclareMathOperator{\Tor}{Tor}
\DeclareMathOperator{\Spec}{Spec}
\DeclareMathOperator{\supp}{supp}
\DeclareMathOperator{\im}{im}
\DeclareMathOperator{\id}{id}
\DeclareMathOperator{\Tr}{Tr}
\DeclareMathOperator{\rank}{rank}
\newcommand{\R}{\mathbb R}
\newcommand{\C}{\mathbb C}
\newcommand{\Z}{\mathbb Z}
\newcommand{\N}{\mathbb N}
\newcommand{\Q}{\mathbb Q}
\begin{document}
\section*{021\quad 张量幂上的Schur--Weyl对偶与双中心化子}
\subsection*{来源与收录范围}
Pavel Etingof, Oleg Golberg, Sebastian Hensel, Tiankai Liu, Alex Schwendner, Dmitry Vaintrob, Elena Yudovina, \emph{Introduction to Representation Theory}, MIT 18.712, Fall 2010, Chapter 4, Sections 4.18--4.19。主目标为 Theorem 4.57 的 Schur--Weyl 对偶；收录 Theorems 4.54--4.55、Lemma 4.56 及 Proposition 4.58 的作者证明，分章 PDF 第18--19页。群版本与 Lie 代数版本合计一个问题。获取日期 2026-09-28。
\par\url{https://ocw.mit.edu/courses/18-712-introduction-to-representation-theory-fall-2010/84358595a02a73bced2c4e363a5d66f0_MIT18_712F10_ch4.pdf}
\par\textbf{据原文上下文补列的记号说明（非引文）：}在下面的应用中取有限维复向量空间 $V$；$S_n$ 置换 $V^{\otimes n}$ 的张量因子，$\mathfrak{gl}(V)$ 按各因子上的作用之和作用，$H_j(x)=\sum_i x_i^j$ 为幂和对称多项式。$V_\lambda$ 是与分拆 $\lambda$ 对应的 Specht 模。
\subsection*{作者命题与人类证明}

\begin{theorem}[4.54, Double Centralizer Theorem]
Let $A,B$ be two subalgebras of the algebra $\operatorname{End}E$ of endomorphisms of a finite dimensional vector space $E$, such that $A$ is semisimple, and $B=\operatorname{End}_A E$. Then:
\begin{enumerate}[label=(\roman*)]
\item $A=\operatorname{End}_B E$ (i.e., the centralizer of the centralizer of $A$ is $A$);
\item $B$ is semisimple;
\item as a representation of $A\otimes B$, $E$ decomposes as $E=\bigoplus_{i\in I}V_i\otimes W_i$, where $V_i$ are all the irreducible representations of $A$, and $W_i$ are all the irreducible representations of $B$. In particular, we have a natural bijection between irreducible representations of $A$ and $B$.
\end{enumerate}
\end{theorem}
\begin{proof}
Since $A$ is semisimple, we have a natural decomposition $E=\bigoplus_{i\in I}V_i\otimes W_i$, where $W_i:=\Hom_A(V_i,E)$, and $A=\bigoplus_i\operatorname{End}V_i$. Therefore, by Schur's lemma, $B=\operatorname{End}_A(E)$ is naturally identified with $\bigoplus_i\operatorname{End}(W_i)$. This implies all the statements of the theorem.
\end{proof}
We will now apply Theorem 4.54 to the following situation: $E=V^{\otimes n}$, where $V$ is a finite dimensional vector space over a field of characteristic zero, and $A$ is the image of $\C[S_n]$ in $\operatorname{End}E$. Let us now characterize the algebra $B$. Let $\mathfrak{gl}(V)$ be $\operatorname{End}V$ regarded as a Lie algebra with operation $ab-ba$.
\begin{theorem}[4.55]
The algebra $B=\operatorname{End}_A E$ is the image of the universal enveloping algebra $U(\mathfrak{gl}(V))$ under its natural action on $E$. In other words, $B$ is generated by elements of the form
\[
\Delta_n(b):=b\otimes1\otimes\cdots\otimes1+1\otimes b\otimes\cdots\otimes1+\cdots+1\otimes1\otimes\cdots\otimes b,
\quad b\in\mathfrak{gl}(V).
\]
\end{theorem}
\begin{proof}
Clearly, the image of $U(\mathfrak{gl}(V))$ is contained in $B$, so we just need to show that any element of $B$ is contained in the image of $U(\mathfrak{gl}(V))$. By definition, $B=S^n\operatorname{End}V$, so the result follows from part (ii) of the following lemma.
\begin{lemma}[4.56]
Let $k$ be a field of characteristic zero.
\begin{enumerate}[label=(\roman*)]
\item For any finite dimensional vector space $U$ over $k$, the space $S^nU$ is spanned by elements of the form $u\otimes\cdots\otimes u$, $u\in U$.
\item For any algebra $A$ over $k$, the algebra $S^nA$ is generated by elements $\Delta_n(a)$, $a\in A$.
\end{enumerate}
\end{lemma}
\begin{proof}
(i) The space $S^nU$ is an irreducible representation of $\operatorname{GL}(U)$ (Problem 3.19). The subspace spanned by $u\otimes\cdots\otimes u$ is a nonzero subrepresentation, so it must be everything.

(ii) By the fundamental theorem on symmetric functions, there exists a polynomial $P$ with rational coefficients such that $P(H_1(x),\ldots,H_n(x))=x_1\cdots x_n$ (where $x=(x_1,\ldots,x_n)$). Then
\[
P(\Delta_n(a),\Delta_n(a^2),\ldots,\Delta_n(a^n))=a\otimes\cdots\otimes a.
\]
The rest follows from (i).
\end{proof}
\end{proof}
Now, the algebra $A$ is semisimple by Maschke's theorem, so the double centralizer theorem applies, and we get the following result, which goes under the name ``Schur--Weyl duality''.
\begin{theorem}[4.57]
\begin{enumerate}[label=(\roman*)]
\item The image $A$ of $\C[S_n]$ and the image $B$ of $U(\mathfrak{gl}(V))$ in $\operatorname{End}(V^{\otimes n})$ are centralizers of each other.
\item Both $A$ and $B$ are semisimple. In particular, $V^{\otimes n}$ is a semisimple $\mathfrak{gl}(V)$-module.
\item We have a decomposition of $A\otimes B$-modules $V^{\otimes n}=\bigoplus_\lambda V_\lambda\otimes L_\lambda$, where the summation is taken over partitions of $n$, $V_\lambda$ are Specht modules for $S_n$, and $L_\lambda$ are some distinct irreducible representations of $\mathfrak{gl}(V)$ or zero.
\end{enumerate}
\end{theorem}
The Schur--Weyl duality for the Lie algebra $\mathfrak{gl}(V)$ implies a similar statement for the group $\operatorname{GL}(V)$.
\begin{proposition}[4.58]
The image of $\operatorname{GL}(V)$ in $\operatorname{End}(V^{\otimes n})$ spans $B$.
\end{proposition}
\begin{proof}
Denote the span of $g^{\otimes n}$, $g\in\operatorname{GL}(V)$, by $B'$. Let $b\in\operatorname{End}V$ be any element.

We claim that $B'$ contains $b^{\otimes n}$. Indeed, for all values of $t$ but finitely many, $t\cdot\operatorname{Id}+b$ is invertible, so $(t\cdot\operatorname{Id}+b)^{\otimes n}$ belongs to $B'$. This implies that this is true for all $t$, in particular for $t=0$, since $(t\cdot\operatorname{Id}+b)^{\otimes n}$ is a polynomial in $t$.

The rest follows from Lemma 4.56.
\end{proof}
\begin{corollary}[4.59]
As a representation of $S_n\times\operatorname{GL}(V)$, $V^{\otimes n}$ decomposes as $\bigoplus_\lambda V_\lambda\otimes L_\lambda$, where $L_\lambda=\Hom_{S_n}(V_\lambda,V^{\otimes n})$ are distinct irreducible representations of $\operatorname{GL}(V)$ or zero.
\end{corollary}

\subsection*{整理说明（非作者正文）}
完整性范围包括中心化子代数的实际生成证明，而不只引用双中心化子定理。作者的 Theorem 4.57 之前已有完成推导的过渡段，没有独立尾随 proof 环境，保留其原顺序。允许的标准先修为 Maschke 定理、Schur 引理、半单代数的矩阵块结构、特征零下对称幂的纯幂张成性/作者所引 Problem 3.19 的不可约性、以及对称函数基本定理；不把 Schur--Weyl 对偶作为先修。Specht 模分类是命名分解因子的先修，不在本条重证。

作者对 $V$ 的一般措辞与 $\C[S_n]$ 同时出现；本文件开头明确采用其复数域情形，不把证明扩充到任意非代数闭域。第19页实际查看了图像，确认幂和公式、分解和群版本正文；第18页图像接口未取得，按其可读 PDF 文本转录。只规范化算子、希腊字母与张量符号，没有把源文件的字体抽取乱码照抄为数学符号。未取得本地 PDF 原件。
\subsection*{可修改的简短标注（非作者正文）}
$c$：张量幂的联合表示分解。

$a$：置换群作用；Lie 代数对角作用；泛包络代数；中心化子；半单代数；重数空间；对称函数幂和；纯张量幂；多项式插值。

\texttt{cross\_theory\_status = yes}。在同一张量幂上把置换群的交换算子识别为 Lie 代数生成的代数，再通过双中心化子与重数空间建立两类不可约表示的配对。位置：Theorems 4.54--4.57 与 Lemma 4.56(ii)。

全部标注均为非穷尽、可修改初稿；未标注不表示未使用。
\subsection*{许可与修改记录}
原文来自 MIT OpenCourseWare，作者与课程如上。按来源网站标示的 Creative Commons Attribution--NonCommercial--ShareAlike 4.0 许可复用；本转录及整理沿用同一许可。2026-09-28：助手转录、加入来源定位与简短标注；不暗示作者或 MIT 认可本次整理。许可：\url{https://creativecommons.org/licenses/by-nc-sa/4.0/}。
\end{document}
